The product of two nonzero finite sums has least exponent equal to the sum of their least exponents. Only the product of their least terms contributes there, and its coefficient is nonzero in $k$. Thus $k[\Gamma]$ is a domain and $v_0(uv)=v_0(u)+v_0(v)$. Addition can only cancel terms, so $v_0(u+v)\geq\min(v_0(u),v_0(v))$ when the sum is nonzero.

On the fraction field define $v(u/w)=v_0(u)-v_0(w)$. If $u/w=u'/w'$, then $uw'=u'w$, and multiplicativity of $v_0$ makes the values equal. Multiplicativity of $v$ follows immediately. For a nonzero sum, putting fractions over a common denominator and applying the inequality for $v_0$ gives the valuation inequality. Any extension must have this formula, so it is unique. Finally, $v(X^\alpha)=\alpha$ for every $\alpha\in\Gamma$, so the value group is all of $\Gamma$.
